\documentclass{article}

\usepackage{amsmath} % math stuff
\usepackage{amssymb} % math stuff
\usepackage{array} % equations and stuff
\usepackage{bm} % bold math
%\usepackage{booktabs} % extra table rule options
%\usepackage{caption} % suppressed table numbering; incompatible with revtex, and longtable, I think
\usepackage{comment} % comment environment
%\usepackage{enumitem} % customization of enumeration, itemize, and description
\usepackage[T1]{fontenc} % font encoding for special characters, must also use scalable font package
\usepackage[margin=0.8in]{geometry} % paper sizes and margins (but be careful not to mess up pre-defined pages)
\usepackage{graphicx} % for graphics
%\usepackage{helvet} % default font is the helvetica postscript font
\usepackage[utf8]{inputenc} % special characters in tex input
\usepackage{layouts} % print units like widths
\usepackage{lipsum} % lorem ipsum filler text
\usepackage{lmodern} % scalable font?
\usepackage{longtable} % multi-page tables
\usepackage{makecell} % specify line-breaks in table cells
\usepackage{mathrsfs} % math script font
\usepackage{mhchem} % easier chemical formula
\usepackage{microtype} % allows disabling of ligatures
\usepackage{multicol} % multicolumns
%\usepackage{newcent} % new century schoolbook font
\usepackage{nicefrac}
\usepackage{numprint} % print and format (large) numbers
\usepackage{parskip} % removes paragraph indentation, and adjusts paragraph skip, as well as list items
\usepackage{pdfpages} % add pdf files as pages
%\usepackage{setspace} % adjust text spacing and indents
\usepackage{siunitx} % decimal alignment, and degree symbol
\usepackage{subfigure} % divided figures
%\usepackage{tabu} % extra table options
\usepackage{textcomp} % symbols
\usepackage{threeparttablex} % better footnotes with longtable
\usepackage{titling} % title placement
\usepackage{ulem} % strikethrough text
\usepackage{upgreek} % upright Greek letters
%\usepackage{url} % superceded by hyperref
\usepackage{verbatim} % verbatim environment
\usepackage{xcolor} % colors and color boxes
\usepackage{xspace} % commands that don't eat up white space
\usepackage{hyperref} % links and page setup; should always come last

\hypersetup{
 bookmarks=true,
 colorlinks=true,
 citecolor=blue,
 linkcolor=blue,
 urlcolor=blue,
 pdfstartview={XYZ null null 1.0} % default open view is 100%
}

\DisableLigatures[f,t]{encoding = T1} % disable ff, fi, fl, tt ligatures; without options, it also disables -- = endash
\renewcommand{\arraystretch}{1.0} % extra vertical (and horizontal?) space in tables

% define centered, left- and right-aligned columns with specified widths
\newcommand{\PreserveBackslash}[1]{\let\temp=\\#1\let\\=\temp}
\newcolumntype{C}[1]{>{\PreserveBackslash\centering}p{#1}}
\newcolumntype{L}[1]{>{\PreserveBackslash\raggedright}p{#1}}
\newcolumntype{R}[1]{>{\PreserveBackslash\raggedleft}p{#1}}

\begin{document}

\pagestyle{empty} % don't number pages

% custom title
\begin{center}
{\LARGE Express Riddler}

\vspace{0.15in}

{\Large 29 April 2022}
\end{center}


\section*{Riddle:}

There are many fractions with a numerator of 1 whose decimal expansions don't go on to infinitely many decimal places.
For example, 1/4 is equivalent to the decimal 0.25, and 1/500 is equivalent to 0.002.
However, the decimal expansion of 1/3 is 0.33333\dots, a decimal that never terminates.

If you were to add up all these numbers---fractions with a numerator of 1 whose decimal expansions don't go on forever---what would be the sum?
(Note: Before you ask, let's include the fraction 1/1 in this group.)


\section*{Solution:}

The series in question is the sum of the inverses of all numbers with only 2 or 5 as factors:

\[
1+\frac{1}{2}+\frac{1}{4}+\frac{1}{5}+\frac{1}{8}+\frac{1}{10}+\frac{1}{16}+\frac{1}{20}+\frac{1}{25}\dots
\]

This can be rewritten as the products of two series, separated into the powers of 2 and 5:

\[
\left(1+\frac{1}{2}+\frac{1}{4}+\frac{1}{8}\dots\right)
\left(1+\frac{1}{5}+\frac{1}{25}+\frac{1}{125}\dots\right)
\]

Each of these two series is geometric and has a finite sum because of the result:

\[
\sum_{n=0}^{\infty}r^{n}=\frac{1}{1-r}
\]

when $|r|<1$.
In this case,

\begin{align*}
\sum_{n=0}^{\infty}\left(\frac{1}{2}\right)^{n}=2 \\
\sum_{n=0}^{\infty}\left(\frac{1}{5}\right)^{n}=\frac{5}{4}
\end{align*}

So the series simplifies to $2(\nicefrac{5}{4})$, and the solution is \fcolorbox{red}{white}{\textbf{\nicefrac{5}{2}}}\,.


\end{document}